\section{Gaussian Space --- Unique Games Definitions}
%%% generated by the blueprint pipeline from TCSlib.BooleanAnalysis.GaussianSpace.UniqueGamesDefinitions
\subsection*{Gaussian Space --- Unique Games Definitions: Introduction}

Declarations of the Lean module \texttt{TCSlib.BooleanAnalysis.GaussianSpace.UniqueGamesDefinitions}.

\subsection*{Gaussian Space --- Unique Games Definitions: Results}

\begin{definition}[Unique game constraint data]
\lean{BooleanAnalysis.GaussianSpace.UniqueGames.UniqueGameConstraint}
\label{BooleanAnalysis.GaussianSpace.UniqueGames.UniqueGameConstraint}
\leanok
A unique game constraint consists of two natural-number variable indices and a list of
natural numbers serving as its table. No further conditions on these data are imposed.
\end{definition}

\begin{definition}[Validity of a unique game constraint]
\lean{BooleanAnalysis.GaussianSpace.UniqueGames.UniqueGameConstraint.Valid}
\label{BooleanAnalysis.GaussianSpace.UniqueGames.UniqueGameConstraint.Valid}
\leanok
\uses{BooleanAnalysis.GaussianSpace.UniqueGames.UniqueGameConstraint}
For a natural number $q$ and a unique game constraint $e$, call $e$ valid if its two
variable indices are distinct and its table is a permutation of the list
$[0,1,\ldots,q-1]$.
\end{definition}

\begin{definition}[Encoding rows of natural numbers]
\lean{BooleanAnalysis.GaussianSpace.UniqueGames.encodeNatRows}
\label{BooleanAnalysis.GaussianSpace.UniqueGames.encodeNatRows}
\leanok
\uses{Std.Sat.CNF.serialize}
Encode a list of rows of natural numbers as a Boolean list by representing each row
entry $k$ as the positive literal with variable index $k$, then serializing the
resulting list of clauses in conjunctive normal form.
\end{definition}

\begin{definition}[Parsing rows of natural numbers]
\lean{BooleanAnalysis.GaussianSpace.UniqueGames.parseNatRows}
\label{BooleanAnalysis.GaussianSpace.UniqueGames.parseNatRows}
\leanok
\uses{Std.Sat.CNF.parse}
Given a list of Boolean values, parse it as a conjunctive-normal-form formula. If
parsing succeeds and every literal in every clause is positive, return the formula as a
list of rows of natural-number variable indices; otherwise return no result.
\end{definition}

\begin{definition}[Parsing a unique game constraint]
\lean{BooleanAnalysis.GaussianSpace.UniqueGames.parseUniqueGameConstraint}
\label{BooleanAnalysis.GaussianSpace.UniqueGames.parseUniqueGameConstraint}
\leanok
\uses{BooleanAnalysis.GaussianSpace.UniqueGames.UniqueGameConstraint, BooleanAnalysis.GaussianSpace.UniqueGames.UniqueGameConstraint.Valid}
For a natural number $q$ and a list of natural numbers, return a unique game constraint
if the list has the form $u,v$ followed by a table that makes the constraint with
indices $u$ and $v$ valid for $q$. Return no result if the list has fewer than two
entries or the resulting constraint is invalid.
\end{definition}

\begin{definition}[Parsing a unique game]
\lean{BooleanAnalysis.GaussianSpace.UniqueGames.parseUniqueGame}
\label{BooleanAnalysis.GaussianSpace.UniqueGames.parseUniqueGame}
\leanok
\uses{BooleanAnalysis.GaussianSpace.UniqueGames.UniqueGameConstraint, BooleanAnalysis.GaussianSpace.UniqueGames.parseNatRows, BooleanAnalysis.GaussianSpace.UniqueGames.parseUniqueGameConstraint}
For a natural number $q$ and a Boolean list, first parse the list as rows of natural
numbers, requiring every literal to be positive. Then parse each row as a valid unique
game constraint for $q$; return the list of constraints if all rows parse successfully,
and no result otherwise.
\end{definition}

\begin{definition}[Encoding unique game constraints]
\lean{BooleanAnalysis.GaussianSpace.UniqueGames.encodeUniqueGame}
\label{BooleanAnalysis.GaussianSpace.UniqueGames.encodeUniqueGame}
\leanok
\uses{BooleanAnalysis.GaussianSpace.UniqueGames.UniqueGameConstraint, BooleanAnalysis.GaussianSpace.UniqueGames.encodeNatRows}
Given a list $G$ of unique game constraints, encode each constraint as the
natural-number row consisting, in order, of its left variable index, its right variable
index, and its table entries; the result is the encoding of the resulting list of rows
as a Boolean list.
\end{definition}

\begin{definition}[Parsing a max-cut edge]
\lean{BooleanAnalysis.GaussianSpace.UniqueGames.parseMaxCutEdge}
\label{BooleanAnalysis.GaussianSpace.UniqueGames.parseMaxCutEdge}
\leanok
For a list of natural numbers, return the ordered pair $(u,v)$ exactly when the list
consists of two entries $u,v$ with $u\ne v$; in every other case, return no pair.
\end{definition}

\begin{definition}[Parsing max-cut edges]
\lean{BooleanAnalysis.GaussianSpace.UniqueGames.parseMaxCut}
\label{BooleanAnalysis.GaussianSpace.UniqueGames.parseMaxCut}
\leanok
\uses{BooleanAnalysis.GaussianSpace.UniqueGames.parseMaxCutEdge, BooleanAnalysis.GaussianSpace.UniqueGames.parseNatRows}
Decode a Boolean list as a list of natural-number rows, and parse each row as an edge by
the rule that it must consist of two distinct natural numbers. Return the resulting list
of ordered pairs if every row is a valid edge, and return no result otherwise.
\end{definition}

\begin{definition}[Encoding max-cut edges]
\lean{BooleanAnalysis.GaussianSpace.UniqueGames.encodeMaxCut}
\label{BooleanAnalysis.GaussianSpace.UniqueGames.encodeMaxCut}
\leanok
\uses{BooleanAnalysis.GaussianSpace.UniqueGames.encodeNatRows}
Given a list $E$ of ordered pairs of natural numbers, encode it as a list of
natural-number rows, using the two entries of each pair as a row, and return the Boolean
encoding of those rows.
\end{definition}

\begin{definition}[Unique game assignment value]
\lean{BooleanAnalysis.GaussianSpace.UniqueGames.uniqueGameAssignmentValue}
\label{BooleanAnalysis.GaussianSpace.UniqueGames.uniqueGameAssignmentValue}
\leanok
\uses{BooleanAnalysis.GaussianSpace.UniqueGames.UniqueGameConstraint}
Let $q$ be a natural number, let $G$ be a list of unique game constraints, and let
$\ell$ assign a label in $\{0,\ldots,q-1\}$ to each natural-number variable. The
assignment value is the number of constraints in $G$ whose table, at the position given
by the label of the right variable, contains the label of the left variable, divided by
the length of $G$; a table position outside its range does not satisfy the constraint,
and when $G$ is empty the value is $0$.
\end{definition}

\begin{definition}[Unique game optimum]
\lean{BooleanAnalysis.GaussianSpace.UniqueGames.uniqueGameOptimum}
\label{BooleanAnalysis.GaussianSpace.UniqueGames.uniqueGameOptimum}
\leanok
\uses{BooleanAnalysis.GaussianSpace.UniqueGames.UniqueGameConstraint, BooleanAnalysis.GaussianSpace.UniqueGames.uniqueGameAssignmentValue}
For a natural number $q$ and a list $G$ of unique game constraints, the unique game
optimum is the supremum of the assignment values over all labelings of the
natural-number variables by the $q$ labels.
\end{definition}

\begin{definition}[Max-cut vertex bound]
\lean{BooleanAnalysis.GaussianSpace.UniqueGames.maxCutVertexBound}
\label{BooleanAnalysis.GaussianSpace.UniqueGames.maxCutVertexBound}
\leanok
\uses{Std.Sat.CNF.numVars}
Given a list $E$ of ordered pairs of natural numbers, form one clause for each pair,
containing its two entries as positive literals. The vertex bound is one plus the
largest variable index occurring in these clauses, or $0$ if no variable occurs.
\end{definition}

\begin{definition}[Weighted graph associated to max-cut edges]
\lean{BooleanAnalysis.GaussianSpace.UniqueGames.maxCutGraph}
\label{BooleanAnalysis.GaussianSpace.UniqueGames.maxCutGraph}
\leanok
\uses{BooleanAnalysis.GaussianSpace.UniqueGames.maxCutVertexBound, BooleanAnalysis.Hypercontractivity.WeightedDigraph}
Given a list $E$ of ordered pairs of natural numbers, let $N$ be one plus the largest
endpoint appearing in $E$, or $0$ if $E$ is empty, and take the vertices to be
$0,\ldots,N-1$. There is a directed edge from $v$ to $w$ exactly when $(v,w)$ or $(w,v)$
occurs in $E$. The weight of the ordered pair $(v,w)$ is the total multiplicity of
$(v,w)$ and $(w,v)$ in $E$, divided by twice the length of $E$; when $E$ is empty, this
weight is $0$.
\end{definition}

\begin{definition}[Yes instances of a unique game]
\lean{BooleanAnalysis.GaussianSpace.UniqueGames.uniqueGameYes}
\label{BooleanAnalysis.GaussianSpace.UniqueGames.uniqueGameYes}
\leanok
\uses{BooleanAnalysis.GaussianSpace.UniqueGames.parseUniqueGame, BooleanAnalysis.GaussianSpace.UniqueGames.uniqueGameOptimum}
For a natural number $q$ and a real threshold $\beta$, the yes language consists of
Boolean strings that parse as some nonempty unique game with alphabet size $q$ and whose
optimum is at least $\beta$.
\end{definition}

\begin{definition}[No instances of a unique game]
\lean{BooleanAnalysis.GaussianSpace.UniqueGames.uniqueGameNo}
\label{BooleanAnalysis.GaussianSpace.UniqueGames.uniqueGameNo}
\leanok
\uses{BooleanAnalysis.GaussianSpace.UniqueGames.parseUniqueGame, BooleanAnalysis.GaussianSpace.UniqueGames.uniqueGameOptimum}
For a natural number $q$ and a real threshold $\alpha$, the no language consists of
Boolean strings that parse as some nonempty unique game with alphabet size $q$ and whose
optimum is strictly less than $\alpha$.
\end{definition}

\begin{definition}[Yes instances of Max-Cut]
\lean{BooleanAnalysis.GaussianSpace.UniqueGames.maxCutYes}
\label{BooleanAnalysis.GaussianSpace.UniqueGames.maxCutYes}
\leanok
\uses{BooleanAnalysis.GaussianSpace.MaxCut.optimum, BooleanAnalysis.GaussianSpace.UniqueGames.maxCutGraph, BooleanAnalysis.GaussianSpace.UniqueGames.parseMaxCut}
For a real threshold $\beta$, the yes language consists of Boolean strings that parse as
some nonempty edge list defining a graph whose Max-Cut optimum is at least $\beta$.
\end{definition}

\begin{definition}[No instances of Max-Cut]
\lean{BooleanAnalysis.GaussianSpace.UniqueGames.maxCutNo}
\label{BooleanAnalysis.GaussianSpace.UniqueGames.maxCutNo}
\leanok
\uses{BooleanAnalysis.GaussianSpace.MaxCut.optimum, BooleanAnalysis.GaussianSpace.UniqueGames.maxCutGraph, BooleanAnalysis.GaussianSpace.UniqueGames.parseMaxCut}
For a real threshold $\alpha$, the no language consists of Boolean strings that parse as
some nonempty edge list defining a graph whose Max-Cut optimum is strictly less than
$\alpha$.
\end{definition}

\begin{definition}[Polynomial-time promise reduction]
\lean{BooleanAnalysis.GaussianSpace.UniqueGames.PolyTimePromiseReducible}
\label{BooleanAnalysis.GaussianSpace.UniqueGames.PolyTimePromiseReducible}
\leanok
\uses{Complexity.PolyTimeComputable}
Given Boolean-string languages $Y,N,Y',N'$, a polynomial-time promise reduction from
$(Y,N)$ to $(Y',N')$ is a polynomial-time computable function on Boolean strings that
maps every string in $Y$ into $Y'$ and every string in $N$ into $N'$.
\end{definition}

\begin{definition}[Promise NP-hardness]
\lean{BooleanAnalysis.GaussianSpace.UniqueGames.PromiseNPHard}
\label{BooleanAnalysis.GaussianSpace.UniqueGames.PromiseNPHard}
\leanok
\uses{BooleanAnalysis.GaussianSpace.UniqueGames.PolyTimePromiseReducible, Complexity.NP}
A pair of Boolean-string languages $(Y,N)$ is promise NP-hard if $Y$ and $N$ are
disjoint and, for every language $L$ in NP, there is a polynomial-time promise reduction
from $(L,L^c)$ to $(Y,N)$, where $L^c$ is the complement of $L$.
\end{definition}

\begin{definition}[Unique Games Conjecture]
\lean{BooleanAnalysis.GaussianSpace.UniqueGames.UniqueGamesConjecture}
\label{BooleanAnalysis.GaussianSpace.UniqueGames.UniqueGamesConjecture}
\leanok
\uses{BooleanAnalysis.GaussianSpace.UniqueGames.PromiseNPHard, BooleanAnalysis.GaussianSpace.UniqueGames.uniqueGameNo, BooleanAnalysis.GaussianSpace.UniqueGames.uniqueGameYes}
For every real $\varepsilon$ with $0<\varepsilon<1/2$, there is a positive natural
number $q$ such that the promise problem whose yes instances are unique games with
alphabet size $q$ and optimum at least $1-\varepsilon$, and whose no instances are
unique games with alphabet size $q$ and optimum less than $\varepsilon$, is promise
NP-hard.
\end{definition}
